\scanpage{153}
\mspage{27}
% archive sheet number: 152

\unsure{en|---}
\[
C_1=\operatorname{card}\underbrace{\mathrm{Fl}\,M^{\pi_1}}=\operatorname{card}\mathrm{Fl}^{\pi_1}M
\]
(où pour tt $\pi\geq\pi_0$, $M^\pi$ désigne la sous-catégorie \inserted{str.}
pleine de $M$ définie par $\mathrm{Filt}^\pi\deleted{\unsure{s}}$, et $\mathrm{Fl}^\pi M=\mathrm{Fl}\,M^\pi$).

Soit $C_0$ une petite partie de $\mathrm{Fl}\,M$, formée de
flèches coquarrables, et posons
\[
\Phi=\widetilde C_0=\Psi^{*}_{\deleted{\unsure{*}}},\quad\text{où}\quad\Psi=C_{0*}=\Phi_*.
\]
Soit $r$ un ordinal $\geq\operatorname{Sup}(\pi_1,C_1)$, et tel que l'on
ait $C_0\subset\mathrm{Fl}^rM$. On suppose que $C_0$ est \unsure{irrédundant|irrédondant}
(i.e. deux objets de $C_0$ isomorphes sont égaux)
plus généralement, que $C_0$ satisfait la \inserted{très anodine}
condition de non pléthoricité (3.11.) relativement à $r$.
On considère la suite transfinie de foncteurs
$p_\alpha=p_\alpha^{C_0}$ sur $\underline{\mathrm{Fl}}\,M$ (cf. \unsure{n°|N.} 1), \deleted{Alors} \deleted{et on rappelle}
\deleted{\unsure{que ce qui}} donnant lieu aux factorisations
fonctorielles d'une $\underline f\in\underline{\mathrm{Fl}}\,M$, $f:X\text{---}Y$, en
\diagram{\begin{tikzcd}[ampersand replacement=\&,column sep=large]
\text{(3.16)} \& \unsure{x|X} \arrow[r,"i_\alpha(f)"] \& X_\alpha(f) \arrow[r,dash,"p_\alpha(f)","{\unsure{\text{---}|\to}}"'] \& Y .
\end{tikzcd}}

On a alors ce qui suit :

$1^\circ)$ On a (pour mémoire) $i_\alpha(f)\in\Phi$ pour tout $\alpha$, et
de plus, \uline{si $\alpha$ est grand devant $r^{\pi_0}$}, on a aussi
$p_\alpha(f)\in\Psi$. Si $r'$ est un cardinal $>r^{\pi_0}$ (p. ex. le
\deleted{\unsure{ses}} cardinal successeur de $r^{\pi_0}$), on peut prendre
pour $\alpha$ le premier ordinal dont le cardinal soit égal
à $r'$.

$2^\circ)$ \deleted{Supposons $\alpha$ choisi grand devant $r^{\pi_0}$. Alors} Pour
tout cardinal $s$ de la forme
\[
s=c^t\quad\text{avec }t\geq r^{\pi_0},\quad c\geq2
\]
(donc $s=2^t$ si on prend $c\leq t\unsure{\}}$, cf lemme 3.\unsure{2}), $\underline{\mathrm{Fl}}^{\,s}M$
est stable sous les $p_\alpha$ tels que $\operatorname{card}\alpha\leq s$. Si on
prend p. ex. pour $\alpha$ le premier ordinal dont le
cardinal soit $>r^{\pi_0}$, alors $\underline{\mathrm{Fl}}^{\,s}M$ est stable par
$p_\alpha\unsure{{}^{*}}$, qui par $1^\circ)$ réalise donc une factorisation fonctorielle
(3.\unsure{16|15}) par $i_\alpha(f)\in\Phi$, $p_\alpha(f)\in\Psi$.
